% SEGMENT OLP-0684-S01
% Part: proof-theory
% Chapter: proof-search
% Section: tableaux

\documentclass[../../../include/open-logic-section]{subfiles}

\begin{document}

\olfileid{pt}{ps}{tab}

\olsection{அட்டவணை மரங்கள்}

அட்டவணை மர அமைப்புகள், தொடரணிக் கணிதங்களுடன் நெருங்கிய
தொடர்புடைய !!{proof} அமைப்புகள். கையால் தேடுவதற்கும்
மென்பொருளில் செயல்படுத்துவதற்கும் ஏற்ற நிறுவல் தேடலை அவை
வழங்குகின்றன. தொடரணிகளின் மரத்தைக் கட்டுவதற்குப் பதிலாக,
அட்டவணை மரத்தில் !!a{proof} என்பது !!{formula}s இன் மரம்.
ஆனால் இயல்பான நிறுவலைப் போலன்றி, இதன் விதிகள் தொடரணிக்
கணித விதிகளை ஒத்திருக்கின்றன. அட்டவணை மர !!{proof} என்பது,
அடிப்படையில், ஒரு மாதிரியை வெளிப்படையாகத் தேடுவது.
எனவே இவற்றைச் சில நேரங்களில் \emph{பொருளியல்} அட்டவணை
மரங்கள் என்றும் \emph{மெய்மை மரங்கள்} என்றும் அழைக்கிறார்கள்.

% SEGMENT OLP-0684-S02
\emph{குறியிட்ட} அட்டவணை மரம் என்பது குறியிட்ட
!!{formula}s இன் மரம். ஒரு குறியிட்ட வாய்பாடு
$\sFmla{\True}{!A}$ அல்லது $\sFmla{\False}{!A}$
என்ற வடிவில் இருக்கும். மரங்கள் கீழ்நோக்கி வளர்கின்றன.
நாம் !!{prove} செய்ய விரும்புவதை வரையறுக்கும்
!!{formula}s மேலே முதலில் இடப்படுகின்றன. எடுத்துக்காட்டாக,
$!A, !B \Sequent !C, !D$ என்ற தொடரணியை
!!{prove} செய்யத் தொடங்கும்போது,
$\sFmla{\True}{!A}, \sFmla{\True}{!B},
\sFmla{\False}{!C}, \sFmla{\False}{!D}$
என்று தொடங்குவோம். $!A$ உம் $!B$ உம் மெய்யாகவும்,
$!C$ உம் $!D$ உம் பொய்யாகவும் அமையும்
!!^a{structure} கிடைத்தால், அந்த !!a{structure},
$!A, !B \Sequent !C, !D$ ஏற்புடையதல்ல என்பதைக் காட்டும்.

% SEGMENT OLP-0684-S03
அட்டவணை மரத்தின் இலைக் கணுவை, அதாவது மிகக் கீழுள்ள
கணுவை, கிளை விரிவாக்க விதிகளால் நீட்டிக்கலாம். அவை அதே
கிளையில் மேலும் கணுக்களைச் சேர்க்கலாம்; அல்லது அதன் கீழ்
இரண்டு உடன்பிறப்புக் கணுக்களைச் சேர்த்து அந்தக் கிளையை
இரு கிளைகளாகப் பிரிக்கலாம். நீட்டிக்கப்படும் இலையின்
மேலே அதே கிளையில் உள்ள எந்தக் குறியிட்ட !!{formula}
க்கும் விதியைப் பயன்படுத்தலாம். மரபுவழி அட்டவணை மர
அமைப்பான \Log{Tc} இன் கிளை விரிவாக்க விதிகள்
\olref{tab:Tc} இல் தரப்பட்டுள்ளன.

\subfile{rules-Tc}

இவ்விதிகள் \Log{G3c} இன் விதிகளைத் தலைகீழாகப்
பயன்படுத்துவதையே காணலாம். $\True$, $\False$ என்ற
குறிகள், ஒரு !!a{formula} முறையே தொடரணியின் இடப்புறத்திலா
வலப்புறத்திலா முதன்மை வாய்பாடாக உள்ளது என்பதைக் காட்டுகின்றன.

% SEGMENT OLP-0684-S04
ஒரு கிளையில் $\sFmla{\True}{!A}$ உம்
$\sFmla{\False}{!A}$ உம் இரண்டும் இருந்தால்,
அக்கிளை \emph{மூடியது}. எல்லாக் கிளைகளும் மூடியிருந்தால்,
முழு அட்டவணை மரமும் மூடியது என்போம். எடுத்துக்காட்டாக,
$!A \lor (!B \land !C) \Sequent (!A \lor !B) \land
(!A \lor !C)$ க்கான மூடிய அட்டவணை மரம் பின்வருகிறது.
இது $\sFmla{\True}{!A \lor (!B \land !C)}$ மற்றும்
$\sFmla{\False}{(!A \lor !B) \land (!A \lor !C)}$
ஆகியவற்றுடன் தொடங்குகிறது:
\begin{oltableau}
  [\sFmla{\True}{\formula{A} \lor (\formula{B} \land \formula{C})},
    just=\TAss, checked
    [\sFmla{\False}{(\formula{A} \lor \formula{B}) \land
        (\formula{A} \lor \formula{C})}, just=\TAss, checked
      [\sFmla{\True}{\formula{A}}, just={\TRule{\True}{\lor}[1]}
        [\sFmla{\False}{\formula{A} \lor \formula{B}},
          just={\TRule{\False}{\land}[2]}, checked
          [\sFmla{\False}{\formula{A}},
            just={\TRule{\False}{\lor}[4]}
            [\sFmla{\False}{\formula{B}},
              just={\TRule{\False}{\lor}[4]}, close]
          ]
        ]
        [\sFmla{\False}{\formula{A} \lor \formula{C}},
          just={\TRule{\False}{\land}[2]}, checked
          [\sFmla{\False}{\formula{A}},
            just={\TRule{\False}{\lor}[4]}
            [\sFmla{\False}{\formula{C}},
              just={\TRule{\False}{\lor}[4]}, close]
          ]
        ]
      ]
      [\sFmla{\True}{\formula{B} \land \formula{C}},
        just={\TRule{\True}{\lor}[1]}, checked
        [\sFmla{\False}{\formula{A} \lor \formula{B}},
          just={\TRule{\False}{\land}[2]}, checked
          [\sFmla{\True}{\formula{B}},
            just={\TRule{\True}{\land}[3]}, move by=2
            [\sFmla{\True}{\formula{C}},
              just={\TRule{\True}{\land}[3]}
              [\sFmla{\False}{\formula{A}},
                just={\TRule{\False}{\lor}[4]}
                [\sFmla{\False}{\formula{B}},
                  just={\TRule{\False}{\lor}[4]},close
                ]
              ]
            ]
          ]
        ]
        [\sFmla{\False}{\formula{A} \lor \formula{C}},
          just={\TRule{\False}{\land}[2]}, checked
          [\sFmla{\True}{\formula{B}},
            just={\TRule{\True}{\land}[3]}, move by=2
              [\sFmla{\True}{\formula{C}},
                just={\TRule{\True}{\land}[3]}
                [\sFmla{\False}{\formula{A}},
                  just={\TRule{\False}{\lor}[4]}
                  [\sFmla{\False}{\formula{C}},
                  just={\TRule{\False}{\lor}[4]},close
                ]
              ]
            ]
          ]
        ]
      ]
    ]
  ]
\end{oltableau}

% SEGMENT OLP-0684-S05
சரிக்குறிகள், அவற்றுக்குரிய !!a{formula} க்கான விதி அதன்
கீழுள்ள எல்லாக் கிளைகளிலும் பயன்படுத்தப்பட்டுவிட்டதைக்
காட்டுகின்றன. $\otimes$ குறி கிளை மூடப்பட்டதைக் காட்டுகிறது.
விதிகளுக்குப் பின்னுள்ள வரி எண்கள், அவை எந்தக் குறியிட்ட
!!{formula}s க்குப் பயன்படுத்தப்பட்டன என்பதைக் காட்டுகின்றன;
அதாவது, குறிப்பிடப்பட்ட வரியில் அதே கிளையில் உள்ள
!!{formula} க்குப் பயன்படுத்தப்பட்டன.

அட்டவணை மரங்களில் நிறுவல் தேடல் எளிது. கட்டம் கட்டமாக
ஒரு மரத்தை உருவாக்குவோம். $0$ ஆம் கட்டத்தில் தொடக்க
!!{formula}s ஐ மேலே எழுதுகிறோம். $n+1$ ஆம் கட்டத்தில்
ஒவ்வொரு இலைக் கணுவுக்கும் பின்வருமாறு செய்கிறோம்:
அதே கிளையில் மேலிருந்து முதலாவதாக உள்ள, இன்னும்
சரிக்குறியிடப்படாத, விரிவாக்க விதி பொருந்தும் குறியிட்ட
வாய்பாட்டைக் கண்டறிந்து அதற்கு விதியைப் பயன்படுத்துகிறோம்.
அந்தக் குறியிட்ட வாய்பாடு இடம்பெறும் ஒவ்வொரு கிளையிலும்
விதியைப் பயன்படுத்திய பின் அதற்குச் சரிக்குறியிடுகிறோம்.
மேலுள்ள எடுத்துக்காட்டு இப் படிமுறையால் உருவாக்கப்பட்டது.

% SEGMENT OLP-0684-S06
$\sFmla{\True}{\lforall}$ மற்றும்
$\sFmla{\False}{\lexists}$ என்ற வடிவங்களைச் சேர்ந்த
குறியிட்ட !!{formula}s ஒருபோதும் சரிக்குறியிடப்படுவதில்லை.
எனவே அவற்றைத் தனியாகக் கையாள வேண்டும். அவை இருந்தால்,
பொருந்தக்கூடிய எல்லாச் சொற்களுக்கும்
$\TRule{\True}{\forall}$ மற்றும்
$\TRule{\False}{\lexists}$ விதிகள் இறுதியில்
பயன்படுத்தப்படுவதை உறுதிசெய்ய வேண்டும். எடுத்துக்காட்டாக,
ஒவ்வொரு கட்டத்திலும், இவ்வடிவங்களைச் சாராத மேலிருந்து
முதலாவது குறியிட்ட !!{formula} ஐக் கையாண்ட பிறகு,
ஒவ்வொரு கிளையிலும் உள்ள இவ்வகை வாய்பாடுகள் அனைத்துக்கும்
அவ்விதிகளைப் பயன்படுத்தலாம். பயன்படுத்த வேண்டிய சொல்~$t$:
கிளையில் ஏற்கெனவே தோன்றும் !!{constant}s இலிருந்தும்
(!!{constant}s ஏதும் இல்லாவிட்டால் $\Obj c_0$ இலிருந்தும்),
தொடக்க !!{formula}s இல் தோன்றும் !!{function}s இலிருந்தும்
உருவாக்கக்கூடிய சொற்களில், அதற்குரிய குறியிட்ட
!!{formula} ஆன $\sFmla{\True}{!B(t)}$ அல்லது
~$\sFmla{\False}{!B(t)}$ கிளையில் ஏற்கெனவே
தோன்றாதபடி உள்ள முதல் சொல்~$t$ ஆகும்.

% SEGMENT OLP-0684-S07
இந்த நிறுவல் தேடலால் மூடிய அட்டவணை மரம் கிடைக்காவிட்டால்,
விரிவாக்கம் முடிந்த, குறியிடப்படக்கூடிய அணுவல்லாத
!!{formula}s அனைத்தும் சரிக்குறியிடப்பட்ட முடிவுள்ள
திறந்த கிளை கிடைக்கும்; அல்லது தேடல் முடிவிலாத,
முடிவுள்ள கிளைப்பிரிப்புகளைக் கொண்ட மரத்தை உருவாக்கும்.
அப்படியான மரத்தில் ஒரு முடிவிலாத கிளை இருக்க வேண்டும்.
\olref[cpl]{sec} இல் செய்ததைப் போலவே, அந்தக் கிளையில்
$\sFmla{\True}{!A}$ இருந்தால் $\Sat{M}{!A}$ உம்,
$\sFmla{\False}{!A}$ இருந்தால் $\Sat/{M}{!A}$ உம்
நிறைவேறும் வகையில் !!a{structure}~$\Struct{M}$ ஐ
வரையறுக்கலாம். இதற்கு $\Theta =
\Setabs{!A}{\sFmla{\True}{!A} \text{ கிளையில் உள்ளது}}$
என்றும் $\Xi =
\Setabs{!A}{\sFmla{\False}{!A} \text{ கிளையில் உள்ளது}}$
என்றும் எடுத்துக்கொள்கிறோம்.

% SEGMENT OLP-0684-S08
இந்த நிறுவல் தேடல் முறையால் கிடைக்கும் மூடிய அட்டவணை
மரங்களை, ஒவ்வொரு கணுவுக்கும் ஒரு தொடரணியை வழங்குவதன்
மூலம் \Log{G3c} நிறுவல்களாக எளிதில் மாற்றலாம். தொடக்க
வாய்பாடுகள் $\Gamma \Sequent \Delta$ என்ற தொடரணிக்குரியவை
என்றால், அவை ஒவ்வொன்றுக்கும் $\Gamma \Sequent \Delta$
என்ற அடையாளத்தை வழங்குகிறோம். ஒரு கிளை,
$\sFmla{\True}{!A}$ என்ற குறியிட்ட !!{formula} க்குக்
கிளை விரிவாக்க விதியைப் பயன்படுத்துவதால் நீட்டிக்கப்பட்டால்,
அதன் இலைக்கு $!A, \Pi \Sequent \Lambda$ என்ற தொடரணியை
வழங்குகிறோம். குறியிட்ட !!{formula} ஆன
$\sFmla{\False}{!A}$ க்குப் பயன்படுத்தினால்,
அந்த இலைக்கு $\Pi \Sequent \Lambda, !A$ என்பதை
வழங்குகிறோம். அட்டவணை மரத்தில்
\TRule{\True}{\star} விதி பயன்படுத்தப்படும் இடத்தில்,
தொடரணிக்கு \LeftR{\star} விதியைப் பயன்படுத்துகிறோம்;
இதில் $!A$ முதன்மை வாய்பாடு. \TRule{\False}{\star}
பயன்படுத்தப்படும் இடத்தில் \RightR{\star} விதியைப்
பயன்படுத்துகிறோம். அந்த விதியின் முன்கோள்கள், மரத்தில்
சேர்க்கப்பட்ட தொடர்புடைய கணுக்களுக்கு வழங்கிய
தொடரணிகளாகும். $\sFmla{\True}{!A}$ உம்
$\sFmla{\False}{!A}$ உம் கொண்டு ஒரு கிளை மூடும்போது,
அதன் இறுதி இலைக் கணுவின் அடையாளம்
$!A, \Pi \Sequent \Lambda, !A$ என்ற வடிவத் தொடரணியாக
இருக்கும். அட்டவணை மரத்தின் அடுத்தடுத்த சில கணுக்களுக்கு
ஒரே தொடரணி வழங்கப்படும்; அவ்வாறு மீண்டும் வரும்
நிகழ்வுகளை நீக்குகிறோம். எடுத்துக்காட்டாக, மேலுள்ள
அட்டவணை மரம் பின்வரும் !!{proof} ஆக மாறுகிறது:
\[\small
\Axiom$!A \fCenter !A, !B$
\RightLabel{\RightR{\lor}}
\UnaryInf$!A \fCenter !A \lor !B$
\Axiom$!A \fCenter !A, !C$
\RightLabel{\RightR{\lor}}
\UnaryInf$!A \fCenter !A \lor !C$
\RightLabel{\RightR{\land}}
\BinaryInf$!A \fCenter (!A \lor !B) \land (!A \lor !C)$
\Axiom$!B, !C \fCenter !A, !B$
\RightLabel{\RightR{\lor}}
\UnaryInf$!B, !C \fCenter !A \lor !B$
\RightLabel{\LeftR{\land}}
\UnaryInf$!B \land !C \fCenter !A \lor !B$
\Axiom$!B, !C \fCenter !A, !C$
\RightLabel{\RightR{\lor}}
\UnaryInf$!B, !C \fCenter !A \lor !C$
\RightLabel{\LeftR{\land}}
\UnaryInf$!B \land !C \fCenter !A \lor !C$
\RightLabel{\RightR{\land}}
\BinaryInf$!B \land !C \fCenter (!A \lor !B) \land
(!A \lor !C)$
\RightLabel{\LeftR{\lor}}
\BinaryInf$!A \lor (!B \land !C) \fCenter (!A \lor !B) \land
(!A \lor !C)$
\DisplayProof
\]

\end{document}
